% Complete Section 8.2, original PDF 108-109 / printed 98-99.
% block: variable-coefficients-heading page:108 type:section
\subsection{\BT{Variable coefficients}{变系数}}\label{sec:variable-coefficients}

% block: weighted-form-inner-product pages:108,109 type:ordinary-paragraph
\BT{Next we begin to move beyond the Hodge Laplacian. A simple variation is to change the inner product on one or more of the base Hilbert spaces $W^k=L^2\Lambda^k$ from the standard $L^2$ inner product to a weighted inner product:}{下面开始超出 Hodge Laplacian 的范围. 一个简单的变化, 是将一个或多个底层 Hilbert 空间 $W^k=L^2\Lambda^k$ 上的内积, 从标准 $L^2$ 内积改为加权内积:}
\[
\langle\omega,\mu\rangle=\int_\Omega\langle A\omega,\mu\rangle\operatorname{vol},\qquad\omega,\mu\in L^2\Lambda^k(\Omega).
\]
\BT{Here $A$ is a function on $\Omega$ for which $A(x)$ is a symmetric positive definite operator on $\operatorname{Alt}^kT_x\Omega$ for each $x\in\Omega$ (e.g., if $\Omega$ is a domain in $\mathbb R^n$, $A$ can be viewed as a function from $\Omega$ to the space of symmetric positive definite matrices of size $\binom nk$). Assuming an upper bound on $A$ and a positive lower bound on its least eigenvalue, the weighted inner product on $W^k$ is equivalent to the usual one, and so the new Hilbert complex inherits all essential properties as the old, such as having closed ranges. The abstract Hodge Laplacian of the complex, however, is different. Leaving the general case to the reader, consider the simple example of the final segment of the de Rham complex on a domain in $\mathbb R^n$. We take as the base spaces the weighted space $L^2\Lambda^{n-1}(\Omega,A)$ described above, where $A$ is an $n\times n$ symmetric positive definite matrix-valued function, and leave the rest of the complex untouched, so the relevant segment is}{这里 $A$ 是 $\Omega$ 上的函数, 对每个 $x\in\Omega$, $A(x)$ 都是 $\operatorname{Alt}^kT_x\Omega$ 上的对称正定算子 (例如, 若 $\Omega$ 是 $\mathbb R^n$ 中的区域, 则 $A$ 可视为从 $\Omega$ 到阶数为 $\binom nk$ 的对称正定矩阵空间的函数). 假定 $A$ 有上界, 且其最小特征值有正下界, 则 $W^k$ 上的加权内积与通常内积等价, 因而新的 Hilbert 复形继承了原复形的所有本质性质, 例如像是闭的. 然而, 该复形的抽象 Hodge Laplacian 有所不同. 将一般情形留给读者, 考虑 $\mathbb R^n$ 中区域上 de Rham 复形最后一个片段这一简单例子. 取上述加权空间 $L^2\Lambda^{n-1}(\Omega,A)$ 作为底层空间, 其中 $A$ 是 $n\times n$ 对称正定矩阵值函数, 复形其余部分保持不变, 因而相关片段为}
\begin{equation}\label{eq:weighted-divergence-complex}
L^2\Lambda^{n-1}(\Omega,A)\xrightarrow{(\operatorname{div},H(\operatorname{div}))}L^2\Lambda^n(\Omega)\to0.
\end{equation}

% block: variable-coefficient-mixed-problem page:109 type:ordinary-paragraph
\BT{The weak form of the abstract Hodge Laplace problem for this complex then seeks $\sigma\in H(\operatorname{div})$, $u\in L^2$ such that}{这个复形的抽象 Hodge Laplace 问题的弱形式, 就是求 $\sigma\in H(\operatorname{div})$, $u\in L^2$, 使得}
\[
\begin{aligned}
\langle A\sigma,\tau\rangle-\langle u,\operatorname{div}\tau\rangle&=0,&\tau&\in H(\operatorname{div}),\\
\langle\operatorname{div}\sigma,v\rangle&=\langle f,v\rangle,&v&\in L^2.
\end{aligned}
\]
\BT{This is the standard weak formulation corresponding to the problem}{这是下述问题相应的标准弱形式}
\[
A\sigma+\operatorname{grad}u=0,\qquad\operatorname{div}\sigma=f,
\]
\BT{or, equivalently, $-\operatorname{div}C\operatorname{grad}u=f$, $\sigma=-C\operatorname{grad}u$, where $C=A^{-1}$, together with the Dirichlet boundary condition $u=0$ on $\partial\Omega$ which, in this weak formulation, is a natural boundary condition. In short, this variable coefficient partial differential operator is again the abstract Hodge Laplacian of a Hilbert complex, albeit a slightly different Hilbert complex than the $L^2$ de Rham complex. It may be discretized with the same finite element spaces as in the case where $A$ is the identity, and the analysis carries over without additional considerations. Many other elliptic partial differential equations with variable coefficients can be treated in a similar fashion.}{或者等价地, $-\operatorname{div}C\operatorname{grad}u=f$, $\sigma=-C\operatorname{grad}u$, 其中 $C=A^{-1}$, 并附有 $\partial\Omega$ 上的 Dirichlet 边界条件 $u=0$; 在此弱形式中, 它是自然边界条件. 简言之, 这一变系数偏微分算子仍是某个 Hilbert 复形的抽象 Hodge Laplacian, 尽管这个 Hilbert 复形与 $L^2$ de Rham 复形稍有不同. 可以用 $A$ 为恒等映射时的相同有限元空间进行离散化, 分析无需额外考察即可推广. 其他许多变系数椭圆型偏微分方程也可用类似方式处理.}
